\chapter{Diseño hardware}
\label{chap:hardware}

\section{Placa adaptadora: objetivos}

El módulo Raspberry Pi Pico~2 trabaja a \SI{3.3}{\volt} y no puede conectarse directamente al
entorno de la máquina. La placa adaptadora (proyecto KiCad \fichero{tfg\_adaptador}) resuelve esa
adaptación y reúne en un solo circuito:

\begin{itemize}
  \item un transceptor RS-485 semidúplex para el bus Modbus del accionamiento;
  \item cuatro entradas digitales protegidas para finales de carrera y otros contactos;
  \item entradas analógicas de \SIrange{0}{5}{\volt} acondicionadas para el ADC del microcontrolador;
  \item cuatro salidas de potencia de lado alto con diagnóstico;
  \item dos zócalos mikroBUS para placas Click de célula de carga;
  \item bornas de tornillo para todo el cableado de campo y taladros de fijación M3.
\end{itemize}

\section{Revisiones del diseño}

El repositorio contiene dos estados del diseño, que conviene distinguir:

\begin{itemize}
  \item \textbf{Revisión A}: los ficheros \fichero{tfg\_adaptador.kicad\_sch} y
        \fichero{.kicad\_pcb}, creados entre el 12 y el 17 de marzo de 2026. Tienen el esquema, la
        colocación de los componentes y los planos de cobre (figuras~\ref{fig:pcb-2d}
        y~\ref{fig:pcb-3d}).
  \item \textbf{Revisión B}: el esquema de la figura~\ref{fig:esquema-revB}, incorporado al
        repositorio como imagen. Es la asignación de pines que usan todos los programas de prueba
        de la placa y el firmware final \verificar{añadir al repositorio los ficheros KiCad de la
        revisión B, que no están versionados, y regenerar las figuras de la PCB}.
\end{itemize}

El análisis de la lista de conexiones de la revisión A, extraída con \codigo{kicad-cli}, muestra una
serie de problemas que la revisión B corrige (tabla~\ref{tab:erratas}). Documentarlos es útil,
porque son errores típicos de un primer diseño.

\begin{longtable}{@{}L{0.18\textwidth}L{0.35\textwidth}L{0.34\textwidth}@{}}
\caption{Problemas detectados en la revisión A y corrección en la revisión B.}\label{tab:erratas}\\
\toprule
\textbf{Bloque} & \textbf{Revisión A (lista de conexiones)} & \textbf{Revisión B} \\
\midrule
\endfirsthead
\toprule
\textbf{Bloque} & \textbf{Revisión A} & \textbf{Revisión B} \\
\midrule
\endhead
RS-485 & La salida RO del transceptor está en la misma red que \SI{3.3}{\volt} y GPIO13. DI, DE y $\overline{\text{RE}}$ están unidas entre sí y a GPIO14 y GPIO15. El transceptor no puede funcionar en semidúplex. & DE y $\overline{\text{RE}}$ a GPIO12 con \emph{pull-down} de \SI{10}{\kilo\ohm}; RO a GPIO13 con \emph{pull-up} de \SI{10}{\kilo\ohm}; DI a GPIO14. \\
Entradas analógicas & Las salidas de los seguidores llegan a GPIO16--19, que no tienen ADC. & Salidas a GPIO26--28 (ADC0--2). \\
Amplificadores & Cuatro MCP6001 sin alimentación: pines V+ y V$-$ sin conectar. & Un MCP6294 cuádruple alimentado a \SI{5}{\volt} con desacoplo de \SI{100}{\nano\farad}. \\
Salidas de potencia & El TPS4H000 se alimenta de VBUS (\SI{5}{\volt} del USB); diagnóstico ($\overline{\text{ST}}$), DIAG\_EN, CL y THER sin conectar; entradas directas. & Alimentación desde la borna VS\_IN; $\overline{\text{ST1}}$--$\overline{\text{ST4}}$ unidas a DST con \emph{pull-up}; DIAG\_EN controlado; resistencias de \SI{10}{\kilo\ohm} en serie con las entradas. \\
Células de carga & Cada zócalo en un bus I\textsuperscript{2}C distinto (GPIO0/1 y GPIO2/3), sin control de EN ni lectura de INT. & Ambos zócalos en I\textsuperscript{2}C1 (GPIO2/3), habilitación EN1 (GPIO1) e interrupciones INT1 (GPIO0) e INT2 (GPIO4). \\
Entradas digitales & Llegan a GPIO4--7. & Llegan a GPIO16--19. \\
Rutado & El fichero de PCB no contiene pistas (0 segmentos): sólo colocación, 23 vías y planos. & \completar{indicar cómo se fabricó la placa definitiva: PCB de revisión B, placa perforada o cableado manual} \\
\bottomrule
\end{longtable}

\begin{figure}[p]
\centering
\includegraphics[width=\textwidth]{esquema_revB.png}
\caption{Esquema de la placa adaptadora, revisión B.}
\label{fig:esquema-revB}
\end{figure}

\section{Microcontrolador}

La placa aloja el módulo Raspberry Pi Pico~2 en sus dos filas de 20 pines, de modo que puede
sustituirse sin desoldar. El módulo se alimenta y comunica por USB. De él se toman
\SI{5}{\volt} (VBUS) para los amplificadores y \SI{3.3}{\volt} (salida 3V3 del regulador del
módulo) para el transceptor y las placas Click. La asignación completa de pines está en el
anexo~\ref{anx:pines}.

\section{Interfaz RS-485}

Se usa el transceptor THVD1400 de Texas Instruments \cite{thvd1400}: semidúplex, alimentación de
\SIrange{3}{5.5}{\volt}, hasta \SI{500}{\kilo\bit\per\second}, 1/8 de carga unitaria (hasta 256
nodos), protección ESD de \SI{\pm 16}{\kilo\volt} HBM y \SI{\pm 12}{\kilo\volt} IEC 61000-4-2 por
contacto, y protección ante fallos de bus de \SI{\pm 16}{\volt}. Esta robustez es adecuada para un
cable que convive con los de potencia del motor.

Las señales DE y $\overline{\text{RE}}$ están unidas y se controlan desde GPIO12. Una resistencia
de \SI{10}{\kilo\ohm} a masa deja el transceptor en recepción durante el arranque, cuando el pin aún
no está configurado, para que no ocupe el bus. La salida RO lleva una \emph{pull-up} de
\SI{10}{\kilo\ohm} para que la línea de recepción no quede flotando. El bus sale por la borna de tres
polos J5 (A, B y masa). El firmware activa el emisor antes de cada trama y lo desactiva
\SI{200}{\micro\second} después del último byte (sección~\ref{sec:prob-uart}). No hay terminación en
la placa: el enlace es punto a punto y corto \verificar{longitud del cable y si el accionamiento
incorpora terminación de \SI{120}{\ohm}}.

\section{Entradas digitales}

Cada entrada digital (DI1--DI4) pasa por una resistencia serie de \SI{1}{\kilo\ohm} y un diodo zéner
BZX384 a masa \cite{bzx384} antes de llegar al GPIO (figura~\ref{fig:esquema-revB}). La resistencia
limita la corriente y el zéner recorta las sobretensiones y las descargas que entran por el cable.
Los finales de carrera, de contacto normalmente cerrado, se conectan entre la borna y masa, y el
firmware activa la \emph{pull-up} interna del GPIO.

La tensión que admite la entrada de forma continua está limitada por la potencia de la resistencia.
Con un zéner de tensión $V_Z$ y una resistencia de \SI{0.125}{\watt} (encapsulado 0805), la tensión
de entrada máxima es
\begin{equation}
  V_{\text{in,max}} = V_Z + \sqrt{P_R\,R} \approx V_Z + \sqrt{0{,}125 \cdot 1000} \approx V_Z + \SI{11.2}{\volt}.
\end{equation}
Para contactos secos es más que suficiente. Si se quisieran conectar sensores de \SI{24}{\volt},
habría que subir la resistencia a unos \SI{10}{\kilo\ohm}
\completar{tensión nominal del zéner BZX384 montado}.

\section{Entradas analógicas}

Las entradas AIN1--AIN3 admiten \SIrange{0}{5}{\volt}. Cada canal tiene un divisor resistivo de
\SI{17}{\kilo\ohm} y \SI{33}{\kilo\ohm}, con relación $33/50 = 0{,}66$, que lleva \SI{5}{\volt} a
\SI{3.3}{\volt}, seguido de un amplificador operacional en seguidor de tensión. El seguidor presenta
al ADC una impedancia de fuente baja, necesaria para cargar el condensador de muestreo en cada
conversión sin error, y aísla el divisor del ADC. Se usa el MCP6294 \cite{mcp629x}, cuádruple,
\emph{rail-to-rail} en entrada y salida y alimentado a \SI{5}{\volt}.

Con el ADC de 12 bits del RP2350 y referencia de \SI{3.3}{\volt}, la resolución referida a la
entrada es
\begin{equation}
  \Delta V_{\text{in}} = \frac{\SI{3.3}{\volt}}{4095 \cdot 0{,}66} \approx \SI{1.22}{\milli\volt}.
\end{equation}
El cuarto amplificador se usa como comparador con un umbral de \SI{2.5}{\volt}, obtenido de
\SI{5}{\volt} con un divisor de dos resistencias de \SI{10}{\kilo\ohm}. Su salida se adapta a
\SI{3.3}{\volt} con otro divisor 17/33 y llega a una entrada digital (COMP).

El valor de \SI{17}{\kilo\ohm} no pertenece a la serie E24. Para la fabricación conviene usar
\SI{17.4}{\kilo\ohm} (E96), con relación 0,655, o recalibrar el factor en el firmware.

\section{Salidas de potencia}

Las cuatro salidas usan el interruptor inteligente de lado alto TPS4H000-Q1 \cite{tps4h000}: cuatro
canales, alimentación de \SIrange{3.4}{40}{\volt}, \SI{1000}{\milli\ohm} de resistencia en
conducción, límite de corriente ajustable entre \num{0.1} y \SI{1.3}{\ampere}, detección de
sobrecorriente, cortocircuito y carga abierta, protección térmica y capacidad para cargas
inductivas. Las entradas IN1--IN4 se gobiernan desde GPIO7--10 a través de resistencias de
\SI{10}{\kilo\ohm}. Las salidas de estado, en drenador abierto, se unen en una línea común de fallo
(DST) con \emph{pull-up}, y la habilitación del diagnóstico se controla con DIAG\_EN. El firmware
final no usa todavía estas salidas; su funcionamiento se verificó con el programa
\fichero{pcbtest/test\_digital\_outputs} \completar{uso previsto de las salidas y resistencia de
límite de corriente elegida}.

\section{Célula de carga}

La célula de carga se conecta a una placa \emph{Load Cell 2 Click} de MikroElektronika
\cite{mikroelc2}, basada en el acondicionador ZSC31014 \cite{zsc31014}, que se inserta en el zócalo
mikroBUS MD1. El chip responde en la dirección I\textsuperscript{2}C \codigo{0x28}, admite
alimentación de \SIrange{2.7}{5.5}{\volt} y ofrece ocho ganancias de preamplificador (×1,5 a
×192). Su señal EN, controlada desde GPIO1 con la corriente de salida del pin configurada a
\SI{12}{\milli\ampere}, permite apagarlo y encenderlo para forzar un reinicio y abrir la ventana del
modo comando. Esta posibilidad resultó imprescindible, y también problemática
(sección~\ref{sec:prob-zsc}).

El segundo zócalo (MD2) comparte el bus y se usó para evaluar una \emph{Load Cell 4 Click} con
NAU7802 \cite{nau7802} \completar{modelo, capacidad nominal y sensibilidad (mV/V) de la célula de
carga, y forma de montaje en el cable}.

\section{Alimentación}

La lógica se alimenta del USB del ordenador a través del módulo Pico. Las salidas de potencia toman
su tensión de la borna VS\_IN y el accionamiento tiene su propia fuente
\completar{tensión y potencia de la fuente del accionamiento}. Las masas se unen en la placa, y la
masa del bus RS-485 se lleva también a la borna J5 para dar al transceptor una referencia común con
el accionamiento.

\section{Placa de circuito impreso}

La PCB de la revisión A es de dos capas, mide \qtyproduct{123 x 75}{\milli\metre} y tiene cuatro taladros M3
en las esquinas. Los componentes pasivos son SMD 0805, los diodos SOD-323, los amplificadores
SOT-23-5, el transceptor SOIC-8 y el interruptor de potencia HTSSOP-20 con almohadilla térmica.
Las bornas de campo están agrupadas en el lado izquierdo, el módulo Pico en el lado derecho con el
conector USB hacia el borde, y los dos zócalos mikroBUS arriba a la derecha, de modo que las placas
Click quedan accesibles.

\begin{figure}[tbp]
\centering
\includegraphics[width=0.92\textwidth]{pcb_revA_2d.png}
\caption{PCB de la revisión A: capas de cobre, máscara y serigrafía superiores.}
\label{fig:pcb-2d}
\end{figure}

\begin{figure}[tbp]
\centering
\includegraphics[width=0.85\textwidth]{pcb_revA_3d.png}
\caption{Vista 3D de la PCB de la revisión A generada con KiCad.}
\label{fig:pcb-3d}
\end{figure}

\section{Validación de la placa}
\label{sec:hw-validacion}

Cada bloque se validó con un programa de prueba independiente antes de la integración
(tabla~\ref{tab:pruebas-hw}).

\begin{longtable}{@{}L{0.35\textwidth}L{0.34\textwidth}L{0.19\textwidth}@{}}
\caption{Pruebas de la placa adaptadora.}\label{tab:pruebas-hw}\\
\toprule
\textbf{Programa} & \textbf{Procedimiento y criterio} & \textbf{Resultado} \\
\midrule
\endfirsthead
\toprule
\textbf{Programa} & \textbf{Procedimiento y criterio} & \textbf{Resultado} \\
\midrule
\endhead
\fichero{pcbtest/test\_analog} & Llevar AIN a GND, \SI{3.3}{\volt} y \SI{5}{\volt}; la tensión reconstruida debe ser \num{0}, \num{3.3} y \SI{5.0}{\volt}. & \completar{lecturas} \\
\fichero{pcbtest/test\_din} & Cambiar el estado de DIN1--4 y comprobar que se detecta cada flanco. & \completar{} \\
\fichero{pcbtest/test\_digital\_outputs} & Conmutar SIN1 cada segundo y medir la salida. & \completar{} \\
\fichero{test\_limit\_switches} & Pulsar cada final de carrera: nivel bajo en reposo y alto activado. & Correcto (usado en todos los ensayos) \\
\fichero{test\_rs485\_loopback} & Puente GP14--GP13 sin transceptor: debe recibirse la trama enviada a 9600 y 115200 baudios. & Correcto \\
\fichero{test\_rs485\_tx/rx} & Transmisión y recepción a través del transceptor. & Correcto \\
\fichero{test\_modbus\_scan} & Lectura periódica del registro 0 del esclavo 1 a 115200 baudios. & Correcto (código 0x00) \\
\fichero{test\_loadcell} & Detección del ZSC31014 en 0x28 y lectura del puente con estado. & Correcto \\
\bottomrule
\end{longtable}
